\documentclass[11pt,a4paper]{article}

\usepackage[utf8]{inputenc}
\usepackage[T1]{fontenc}
\usepackage{amsmath,amssymb,amsthm}
\usepackage{graphicx}
\usepackage{booktabs}
\usepackage{hyperref}
\usepackage{xcolor}
\usepackage{tcolorbox}
\usepackage[margin=25mm]{geometry}
\usepackage{xeCJK}
\setCJKmainfont{Hiragino Mincho ProN}
\setCJKsansfont{Hiragino Sans}

\newtheorem{observation}{観察}
\newtheorem{coincidence}{数値的一致}
\newtheorem{conjecture}{予想}

\title{$B_2$ 統一仮説の深掘り:\\
Bernoulli ladder と Spec$(\mathbb{Z})$ alphabet}
\author{Wright Brothers Research Program}
\date{2026年4月}

\begin{document}
\maketitle

\begin{abstract}
先行 note~\cite{b2-note}で提案した $B_2$ 統一仮説を深掘りし，
いくつかの興味深い数値的観察を報告する．主要な新しい発見:
(1) \textbf{Bernoulli ladder} $2n/|B_{2n}| = 12, 120, 252, 240, 132, \ldots$
は $1/|\zeta(1-2n)|$ の数列であり，最初の 3 項の積がちょうど $9!$
に等しい，(2) Standard Model の gauge boson 数 $12$ (電弱対称性
破れ後) が $2/|B_2|$ と一致する，(3) $n = 12$ での Bernoulli ladder
の「破綻」(irregular prime 691 の出現)が物理的 cutoff に対応する
可能性，(4) レプトン世代と $\zeta(1-2n)$ 階層の対応仮説
(a_e: $\zeta(-1)$, a_μ: $\zeta(-5)$ が muon g-2 anomaly $\sim 252$ と
合致),  (5) $\Omega_\Lambda = 8\pi B_2^2$ という新しい簡潔表現．
これらは Spec$(\mathbb{Z})$ 上の arithmetic 不変量が物理定数の
``alphabet'' を成すという仮説を強める証拠群であるが，同時に
look-elsewhere effect のリスクを honest に評価する．
\end{abstract}

\tableofcontents

\section{序: Bernoulli ladder}

\subsection{定義}

Bernoulli 数と Riemann zeta の整数値は以下で関係付けられる:
\begin{equation}
\zeta(1 - 2k) = -\frac{B_{2k}}{2k}, \quad k = 1, 2, 3, \ldots
\end{equation}

絶対値を取ると:
\begin{equation}
\frac{1}{|\zeta(1 - 2k)|} = \frac{2k}{|B_{2k}|}
\end{equation}

これを\textbf{Bernoulli ladder}と呼ぶ．具体的な値:

\begin{center}
\begin{tabular}{lllll}
\toprule
$k$ & $B_{2k}$ & $|B_{2k}|$ & $2k/|B_{2k}|$ & $= 1/|\zeta(1-2k)|$ \\
\midrule
1 & $1/6$ & $1/6$ & $12$ & $1/|\zeta(-1)|$\\
2 & $-1/30$ & $1/30$ & $120$ & $1/|\zeta(-3)|$\\
3 & $1/42$ & $1/42$ & $252$ & $1/|\zeta(-5)|$\\
4 & $-1/30$ & $1/30$ & $240$ & $1/|\zeta(-7)|$\\
5 & $5/66$ & $5/66$ & $132$ & $1/|\zeta(-9)|$\\
6 & $-691/2730$ & $691/2730$ & $32760/691$ & $\neq$ integer\\
\bottomrule
\end{tabular}
\end{center}

\textbf{顕著な事実}: 最初の 5 項 ($k = 1$ から $5$) は全て整数．第 6 項
($k = 6$) から突然\textbf{非整数}に転じる．これは $B_{12}$ の
numerator に\textbf{irregular prime $691$} が現れるためで，
Kummer の FLT 研究以来知られる数論的現象~\cite{Kummer1850}．

\subsection{Wright Brothers プロジェクトでの既出}

このラダーの最初の項は既に Wright Brothers の各所に現れている:

\begin{itemize}
\item $12 = 2/|B_2| = 1/|\zeta(-1)|$: $\alpha^{-1}$ 公式の leading 項，
$\Omega_\Lambda$ 公式の分母 (via $\zeta(-1)$)
\item $120 = 4/|B_4| = 1/|\zeta(-3)|$: $\alpha^{-1}$ 公式の次の項．
また 3D Casimir の Boyer 係数 $7/120$
\item $252 = 6/|B_6| = 1/|\zeta(-5)|$: 先行観察でミューオン $g-2$
との関連で言及~\cite{muonNote}
\item $240, 132$: 未使用だが整数であることが注目される
\end{itemize}

\section{面白い発見 1: $9!$ との一致}

\subsection{観察}

ラダーの最初の 3 項を掛け合わせる:
\begin{equation}
12 \times 120 \times 252 = ?
\end{equation}

計算:
\begin{align}
12 \times 120 &= 1440\\
1440 \times 252 &= 362880
\end{align}

一方:
\begin{equation}
9! = 9 \times 8 \times 7 \times 6 \times 5 \times 4 \times 3 \times 2 \times 1 = 362880
\end{equation}

\begin{coincidence}[$9!$ identity]
\begin{equation}
\boxed{\frac{2}{|B_2|}\cdot\frac{4}{|B_4|}\cdot\frac{6}{|B_6|} = 9!}
\end{equation}
等価的:
\begin{equation}
|\zeta(-1)|\cdot|\zeta(-3)|\cdot|\zeta(-5)| = \frac{1}{9!}
\end{equation}
\end{coincidence}

\subsection{これは自明か?}

偶然の可能性: 3 つの整数を掛けた積が factorial に一致するのは，
何かが偶然に合う範疇ではない．$9! = 362880$ は具体的な値で，
$12 \times 120 \times 252$ も具体的で，両者が一致するには何か構造がある．

実際，この identity は\textbf{既知}である可能性がある．確認が必要だが，
Bernoulli 数と Riemann zeta の性質から数論的に証明できるかもしれない．

試みに $k$ を増やす:
\begin{align}
12 \times 120 &= 2! \times 6! \ ? \quad 2 \times 720 = 1440 \checkmark\\
12 \times 120 \times 252 &= 9!\\
12 \times 120 \times 252 \times 240 &= ?
\end{align}
$362880 \times 240 = 87,091,200$
$10! = 3,628,800$
$11! = 39,916,800$
$12! = 479,001,600$
$87,091,200$ は factorial ではない．

$87,091,200 / 10! = 87091200 / 3628800 = 24 = 4!$

そして $87091200 = 10! \times 4! = 3628800 \times 24$

検証: $10! \times 4! = 10! \times 24$ か? $10! = 3628800$, $\times 24 = 87091200$ ✓

So:
\begin{align}
\prod_{k=1}^{3} \frac{2k}{|B_{2k}|} &= 9! \\
\prod_{k=1}^{4} \frac{2k}{|B_{2k}|} &= 10! \times 4!
\end{align}

第 5 項まで:
$87,091,200 \times 132 = 11,496,038,400$
$11! = 39,916,800$
$12! = 479,001,600$
$11,496,038,400 / 11! = 288$
$288 = 2^5 \times 3^2$, not a simple factorial

Hmm. So the pattern "product = factorial" is only clean for the first 3 terms (= 9!).

\subsection{$9!$ の物理的意味?}

$9!$ が現れる物理コンテクスト:
\begin{itemize}
\item 9 次元 sphere $S^9$ の体積係数
\item SU(3) 表現論 (dim 8 隣接表現は SU(3) 自身)
\item E$_8$ 格子の関連
\item Type I string の anomaly coefficient
\end{itemize}

具体的に $9! = 362880$ を含む物理公式は不明だが，$9$ が $3 \times 3$
(generations × colors?) とも解釈できる.

これは\textbf{coincidence vs structure}の grey zone にある観察．
今後の調査課題.

\section{面白い発見 2: Standard Model gauge boson count}

\subsection{観察}

Standard Model の gauge group は $SU(3) \times SU(2) \times U(1)_Y$．
Generator の数 (= gauge boson の数):
\begin{equation}
\dim[SU(3)] + \dim[SU(2)] + \dim[U(1)] = 8 + 3 + 1 = 12
\end{equation}

電弱対称性破れ後，物理的 gauge boson として 12 個:
\begin{itemize}
\item 8 gluon ($g^a$, $a = 1, \ldots, 8$)
\item $W^+, W^-, Z, \gamma$ (4 個)
\end{itemize}
合計 $8 + 4 = 12$.

そして:
\begin{equation}
12 = \frac{2}{|B_2|} = \frac{1}{|\zeta(-1)|}
\end{equation}

\begin{coincidence}[Gauge boson count]
$$\#(\text{SM gauge bosons}) = 12 = 2/|B_2|$$
\end{coincidence}

\subsection{解釈の試み}

この一致は意味があるか?

\textbf{可能性 A: 偶然}
SM の gauge 構造は群論と phenomenology で決まり，Bernoulli 数とは
独立．12 という数はたまたま一致している．

\textbf{可能性 B: 深い構造}
Bernoulli 数は Spec$(\mathbb{Z})$ 上の基本的な arithmetic 不変量．
Spec$(\mathbb{Z})$ が物理の基盤なら，gauge 群の\textbf{次元そのもの}が
arithmetic 的に制約されていても不思議ではない．$12 = 2/|B_2|$ は
「最小の non-trivial arithmetic 不変量」で，SM の gauge 構造は
これに最適化されている．

\textbf{検証}: $\zeta(-1)$ の arithmetic 情報が SU(3)×SU(2)×U(1) の
選択に帰着する first-principles 論理を構築すべき．Connes, Marcolli の
non-commutative geometry アプローチは近い方向~\cite{ConnesMarcolli}．

\textbf{予測}: 仮説 B が正しければ，他の $2k/|B_{2k}|$ も gauge
theoretic な意味を持つはず．例えば $120 = 4/|B_4|$ は $SO(16)$
($\dim = 120$) に対応? E$_8$ ($\dim = 248$) とは合わないが $SO(16)$ の
$\dim = 120$ は合致する．

$SO(16)$ は E$_8 \times$ E$_8$ heterotic string の gauge 群の一部．

\begin{center}
\begin{tabular}{lll}
\toprule
ladder 値 & 候補群 & 次元 \\
\midrule
12 & SM gauge & 12 \\
120 & $SO(16)$ & 120 \\
252 & ? & 252 \\
240 & E$_8$ root system & 240 \\
132 & ? & 132 \\
\bottomrule
\end{tabular}
\end{center}

$240$ と E$_8$ root system の合致は\textbf{特に興味深い}．E$_8$ の
root は 240 個で，$240 = 8/|B_8|$ と一致．

\begin{coincidence}[E$_8$ roots]
$$\text{\#(E$_8$ roots)} = 240 = \frac{8}{|B_8|}$$
\end{coincidence}

もしこれが偶然でなければ，E$_8$ 統一 (M-theory 等で中心的) は
Bernoulli arithmetic に由来する可能性がある．

\subsection{$252$ は何か}

$252$ に対応する群・構造は:
\begin{itemize}
\item $SO(13)$ の dim $= 78$, 違う
\item $\binom{10}{5} = 252$: 二項係数
\item $E_7$ の adjoint の dim + 1 $= 133 + 119 = 252$? 133 + 119 = 252.
\end{itemize}

$\binom{10}{5} = 252$ は面白い．10 個から 5 個選ぶ組合せ数．

他: muon g-2 の anomaly $\sim 252$ ppm？実際の値は $\Delta a_\mu \sim 251 \times 10^{-11}$ で，\textbf{ちょうど 252}に近い．

\begin{coincidence}[Muon g-2]
ミューオン $g-2$ 異常 $\Delta a_\mu \approx 252 \times 10^{-11}$
(Brookhaven + Fermilab 最新値)
\end{coincidence}

これは先行研究~\cite{muonNote}で既に指摘されていた．新しい洞察は，
$252 = 1/|\zeta(-5)|$ であり，Bernoulli ladder 上の位置付けが
明確になったこと．

\section{面白い発見 3: 世代階層仮説}

\subsection{仮説}

SM には 3 世代のレプトン (e, μ, τ) がある．各世代の anomalous
magnetic moment には階層がある．仮説:

\begin{conjecture}[世代-ζ 対応]
レプトン世代 $n$ ($n = 1$ for $e$, $2$ for $\mu$, $3$ for $\tau$) の
anomalous magnetic moment の主要な非標準寄与は $|\zeta(1-2n)|$ に
比例する:
\begin{equation}
\Delta a_n \propto |\zeta(1-2n)| = \frac{|B_{2n}|}{2n}
\end{equation}
\end{conjecture}

具体的:
\begin{itemize}
\item 電子 ($n=1$): $|\zeta(-1)| = 1/12$
\item ミューオン ($n=2$): $|\zeta(-3)| = 1/120$ ? or $|\zeta(-5)|$?
\item タウ ($n=3$): $|\zeta(-5)| = 1/252$
\end{itemize}

もしくは，世代の index を $n = 0, 1, 2$ として:
\begin{itemize}
\item 電子 ($n=0$): $|\zeta(-1)|$ (standard QED)
\item ミューオン ($n=1$): $|\zeta(-3)|$?
\item タウ ($n=2$): $|\zeta(-5)|$?
\end{itemize}

### 数値チェック

電子 g-2: $a_e = 0.0011596521884 \ldots$. Schwinger $= \alpha/(2\pi) = 0.00116141$ with tiny QED corrections. 非常に精密に SM.

ミューオン g-2: 近年の Fermilab + Brookhaven results give $a_\mu^{\rm exp} - a_\mu^{\rm SM} = (251 \pm 59) \times 10^{-11}$ (from BNL+FNAL combined, pre-2023). 最新 lattice QCD hadronic vacuum polarization (BMW collaboration) は SM との差を小さくする可能性があるが，実験差は persist.

Wright Brothers 予測 (仮説): $\Delta a_\mu \propto 1/|\zeta(-3)|$ or
$1/|\zeta(-5)|$?

$1/|\zeta(-3)| = 120$
$1/|\zeta(-5)| = 252$
Observed: $\sim 251 \times 10^{-11}$

$252$ に\textbf{ほぼ完全一致}. したがって ``muon に対応する $\zeta$ 値''
は $\zeta(-5)$ の可能性が高い．

この対応が正しければ:
\begin{itemize}
\item 電子: $\zeta(-1)$ (世代 1 $\leftrightarrow$ 最小 negative 整数)
\item ミューオン: $\zeta(-5)$ (世代 2 $\leftrightarrow$ 次の値)
\item タウ: $\zeta(-9)$? (世代 3 $\leftrightarrow$ $|\zeta(-9)|= 1/132$)
\end{itemize}

ただし $\zeta(-3)$ と $\zeta(-7)$ がスキップされる理由は不明．

\textbf{代替仮説}: 全ての odd negative integer ($\zeta(-1), \zeta(-3),
\zeta(-5), \ldots$) が使われ，世代は「累積」．電子は $\zeta(-1)$，
ミューオンは $\zeta(-1) + \zeta(-5)$ (e + 追加)，etc.

### タウ g-2 予測

タウ lepton の anomalous g-2 は現在\textbf{未測定}(τ の寿命が短くて
直接測れない)．現在の上限:
\begin{equation}
-0.052 < a_\tau < 0.013 \quad\text{(LEP experiment, 95\% CL)}
\end{equation}
$\sim 10^{-2}$ precision. 標準モデル値は $a_\tau^{\rm SM} \approx 1.177 \times 10^{-3}$．

Wright Brothers 予測 (世代-$\zeta$ 対応): $\Delta a_\tau$ は何らかの
$\zeta$ 値に比例するはず．もし $\zeta(-9)$ なら $\Delta a_\tau \propto 1/132$.

これは現状の実験 sensitivity では検証できないが，将来の Belle II,
LHCb upgrade で改善される可能性がある．

\begin{conjecture}[タウ g-2 予測]
$\Delta a_\tau \sim O(10^{-3})$ で，Bernoulli ladder value $1/132$ または
$1/240$ に比例する非標準寄与を持つ．
\end{conjecture}

\section{面白い発見 4: 691 における断絶}

\subsection{irregular prime の登場}

Bernoulli ladder は $k \leq 5$ で整数値を保つが，$k = 6$ で
\begin{equation}
\frac{12}{|B_{12}|} = \frac{12 \times 2730}{691} = \frac{32760}{691}
\end{equation}
となり\textbf{非整数化}する．$691$ は\textbf{最初の irregular prime}の
一つ (他は $37, 59, 67, \ldots$)．

irregular prime とは: 素数 $p$ が $B_2, B_4, \ldots, B_{p-3}$ の
numerators のどれかを割る，という Kummer の定義．$691$ は $B_{12}$ の
numerator を割る．

\subsection{Wright Brothers での物理的意味}

仮説: Bernoulli ladder の「整数性」が壊れる点は，物理的な
\textbf{cutoff} に対応するのかもしれない．

具体的: もし各 $k$ が「何らかの物理量」に対応するなら (例: 12 = SM
gauge, 120 = SO(16), 240 = E_8, 252 = muon g-2)，$k = 6$ では
対応する物理量が整数ではなく「分数的」になる．これは:

\begin{itemize}
\item 物理的に意味のある整数では表せない
\item 新しい結合定数 or $p$-adic 構造が必要
\item 「Wright Brothers framework」の validity 限界
\end{itemize}

もしかすると $k = 1, 2, 3, 4, 5$ がそれぞれ (SM, SO(16), Muon, E_8, ?)
と対応し，$k = 6$ 以上は\textbf{我々の宇宙には物理的 counterpart が
ない}という解釈．

\subsection{$691$ 自体の物理性}

$691$ という数は物理定数として知られていない．しかし:
\begin{itemize}
\item 近い値: $688 = 2^4 \times 43$, $689 = 13 \times 53$, $691$ 素数, $692 = 2^2 \times 173$
\item 物理単位での $691$: 691 keV, 691 GeV, 691 Hz, 等，何も標準的な
\item 数論的 $691$: Ramanujan's tau function mod 691 の congruence,
partition function, Mathieu group M_{12}, ...
\end{itemize}

Ramanujan の \textbf{$\tau$ function congruence}:
$\tau(n) \equiv \sigma_{11}(n) \pmod{691}$ ここで $\tau$ は
modular discriminant $\Delta$ の coefficient．これは $\Delta$ と
Eisenstein series $E_{12}$ の congruence から来る．

E$_{12}$ は $E_{12} = 1 + (65520/691)\sum \sigma_{11}(n) q^n$ で，
691 が分母に登場．

もし Wright Brothers framework が modular form と接続するなら，
$691$ は物理的にも意味を持つかもしれない．しかし現時点では speculation．

\section{面白い発見 5: $\Omega_\Lambda = 8\pi B_2^2$}

\subsection{簡潔表現}

先行研究で $\Omega_\Lambda = 2\pi/9$ が WDW 時間的 DN 仮説から
導出された．別表現として:
\begin{equation}
\Omega_\Lambda = \frac{2\pi}{9} = \frac{2\pi}{9}
\end{equation}
を $B_2$ で書き直す:
\begin{equation}
\Omega_\Lambda = 8\pi B_2^2 = 8\pi \times \frac{1}{36} = \frac{8\pi}{36} = \frac{2\pi}{9}
\end{equation}

\begin{coincidence}[$B_2^2$ form]
\begin{equation}
\boxed{\Omega_\Lambda = 8\pi B_2^2}
\end{equation}
これは $\Omega_\Lambda$ を $B_2$ の\textbf{二乗}と $\pi$ のみで
表現する最も簡潔な形．
\end{coincidence}

\subsection{なぜ二乗か}

物理的解釈の candidate:
\begin{itemize}
\item Casimir 揺らぎの\textbf{power spectrum}は零点揺らぎの
self-correlation $\langle 0|\phi^2|0\rangle$ で，\textbf{$B_2$ の
二乗}を含む
\item 宇宙の\textbf{2 つの境界}(Big Bang + de Sitter) からの寄与を
独立に掛ける
\item Bogoliubov 変換の$|\beta|^2$ 依存性
\item Entanglement entropy の 2nd Renyi
\end{itemize}

これらは speculative だが，「$B_2$ が単独ではなく $B_2^2$ として
現れる」のは物理的な pairing 構造を示唆する．

\subsection{$\pi$ の起源}

$8\pi$ の factor は Einstein equation $G_{\mu\nu} = 8\pi T_{\mu\nu}$
(in $G=c=1$ units) の $8\pi$．これは球の面積係数由来．

したがって:
\begin{equation}
\Omega_\Lambda = \underbrace{8\pi}_{\text{gravity}}\cdot\underbrace{B_2^2}_{\text{arithmetic}}
\end{equation}
という解釈: 「重力と数論の積」が宇宙論定数を与える．

\section{面白い発見 6: $\pi^2/36$ と $\Omega_{\rm dm}$}

\subsection{Dark matter 試行}

先行 note で試みた $\Omega_{\rm dm} \approx \pi^2/36 = 0.2742$ を
再検討．観測値 $0.2645 \pm 0.0045$ (Planck 2018)．差 3.7\%．

$\pi^2/36 = \pi^2 B_2^2 = B_2^2 \cdot \zeta(2)\cdot 6 = \ldots$

整理: $\pi^2/36 = \pi^2/6^2 = \zeta(2)/6$ (since $\zeta(2) = \pi^2/6$)

よって $\Omega_{\rm dm}$ 試行的表現:
\begin{equation}
\Omega_{\rm dm} \stackrel{?}{=} \frac{\zeta(2)}{6} = B_2 \zeta(2)
\end{equation}

\subsection{Flat 性との両立}

Flat 宇宙では $\Omega_{\rm total} = 1$:
\begin{align}
\Omega_\Lambda + \Omega_{\rm dm} + \Omega_b &= 1\\
\frac{2\pi}{9} + B_2 \zeta(2) + \Omega_b &= 1\\
0.6981 + 0.2742 + \Omega_b &= 1\\
\Omega_b &= 0.0277
\end{align}

観測 baryon density: $\Omega_b^{\rm obs} = 0.0493$．差 $\sim 44\%$，
これは\textbf{大きすぎる}．

したがって $\Omega_{\rm dm} = \pi^2/36$ 仮説は flat 性と両立せず，
単独では不整合．

### 代替: Ω_baryon も arithmetic か?

$\Omega_b = 0.0493 \approx 1/20.3$. Can we express 20.3 arithmetically?

$20.3 \approx 6 \pi$? $6\pi = 18.85$. No.
$\pi^3/2 = 15.5$. No.
$2\pi^2 = 19.74$. Close but 3\% off.

Or: $\Omega_b = 1/(2\pi^2) \approx 0.0507$. Observed 0.0493. 3\% off.

### Alternative flat assumption

$\Omega_\Lambda = 2\pi/9$, $\Omega_{\rm dm} = ?$, $\Omega_b = 1/(2\pi^2)$?
Check: $\Omega_{\rm dm} = 1 - 2\pi/9 - 1/(2\pi^2) = 1 - 0.6981 - 0.0507 = 0.2512$
Observed $\Omega_{\rm dm} = 0.2645$. Difference 5\%.

Not clean.

\begin{observation}[dm は arithmetic では微妙]
Dark matter 密度 $\Omega_{\rm dm}$ は $\pi^2/36$ に 3\% で合うが
flat 性を崩す．純粋に arithmetic な表現は現状見つからない．
dark matter 密度は baryogenesis や freeze-out の動的過程に
依存するので，Wright Brothers framework で自然に出る保証はない．
\end{observation}

\section{統合: Spec$(\mathbb{Z})$ alphabet 仮説}

\subsection{命題}

これまでの観察を総合し，以下の統合仮説を提案する:

\begin{conjecture}[Spec$(\mathbb{Z})$ alphabet]
物理定数は，Spec$(\mathbb{Z})$ 上の arithmetic 不変量
$\{B_{2k}, \zeta(s), p_i\}$ (Bernoulli 数, zeta 特殊値, 素数)を
``alphabet'' として書かれる．異なる物理量は異なる ``word''
(arithmetic 不変量の組合せ) を使うが，共通の alphabet に属する．
\end{conjecture}

\subsection{既出の ``words''}

\begin{center}
\begin{tabular}{lll}
\toprule
物理量 & Word & 観測との一致 \\
\midrule
$\alpha^{-1}$ & $2/|B_2| + 4/|B_4| + \ldots$ & $10^{-7}$ \\
$\Omega_\Lambda$ & $8\pi B_2^2 = 2\pi/9$ & 2\% \\
Muon $g-2$ & $6/|B_6| = 252$ & $\sim$ ppm \\
$\alpha^{-1}(m_Z)$ & $1/|\zeta(-3)| + \dim SU(3) = 128$ & $0.04\%$ \\
SM gauge \#\, & $2/|B_2| = 12$ & exact \\
E$_8$ roots & $8/|B_8| = 240$ & exact \\
\bottomrule
\end{tabular}
\end{center}

\subsection{予測される ``words''}

仮説が正しければ，未予測の constants にも arithmetic 表現があるはず:

\begin{itemize}
\item \textbf{Fine-structure at Planck}: $\alpha^{-1}(M_P)$ は何か?
恐らく $2/|B_2| + 4/|B_4| + 6/|B_6| + \ldots$ の partial sum
\item \textbf{Proton/electron mass ratio}: $m_p/m_e = 1836.15$
→ arithmetic 表現探索
\item \textbf{Higgs self-coupling}: $\lambda_H$
→ ?
\item \textbf{Dark matter mass}: 未測定だが framework で予測?
\item \textbf{CKM mixing angles}: arithmetic で?
\end{itemize}

\subsection{例: 2 つの独立な予測の重ね合わせ}

$\alpha^{-1}$ (電磁) と $\Omega_\Lambda$ (重力) が同じ alphabet を
使うなら，両者の組合せも意味を持つはず．

Attempt:
\begin{equation}
\alpha^{-1} \cdot \Omega_\Lambda \approx 137.036 \times 0.6981 = 95.67
\end{equation}

$95.67$ を arithmetic で表現可能か?
\begin{itemize}
\item $30.44 \pi \approx 95.67$: $30.44$ は何か? $30.44 \approx 2 \cdot 1/|B_4|/4 \cdot ...$. 不明
\item $16 - 0.06 \times \ldots$: 単なるフィッティング
\item $9!/(48\pi) \approx 2406.7$: 大きすぎる
\end{itemize}

Clean な表現は見つからない．ただし:
\begin{equation}
\alpha^{-1} \cdot \Omega_\Lambda \cdot |B_2| = 15.94 \approx 2^4 - 0.06
\end{equation}
$\approx 16 = 2^4$ に 0.3\% 近い．

この「$\alpha^{-1} \cdot \Omega_\Lambda \cdot |B_2| \approx 16$」関係は
\textbf{look-elsewhere effect のリスクが高い}偶然かもしれない．ただし
$16 = 2^4$ は spinor 表現の次元などで意味を持ちうる．

\begin{observation}[speculative]
もし $\alpha^{-1} \cdot \Omega_\Lambda \cdot |B_2| = 2^4$ が厳密に成り立つ
なら:
\begin{equation}
\alpha^{-1} = \frac{2^4}{\Omega_\Lambda \cdot |B_2|} = \frac{16}{2\pi/9 \cdot 1/6}
= \frac{16 \cdot 54}{2\pi} = \frac{864}{2\pi} \approx 137.51
\end{equation}
観測 $137.036$ との差 0.34\%. 合わないが「$2^4$ 仮説」は\textbf{reject
ではない}．
\end{observation}

Clean な統合には至らない．

\section{Honest 評価と look-elsewhere effect}

\subsection{発見の格付け}

報告した観察を evidence level で格付け:

\begin{center}
\begin{tabular}{lll}
\toprule
観察 & 強度 & 偶然確率 \\
\midrule
$12 \cdot 120 \cdot 252 = 9!$ & 中 & 既知の identity の可能性 \\
SM gauge \# $= 12$ & 中 & Group theory 独立 \\
E$_8$ roots $= 240$ & 中 & 興味深い \\
Muon $g-2 \sim 252$ & 強 & 先行観察，精度高 \\
$691$ での断絶 & 弱 & 物理対応未確認 \\
$\Omega_\Lambda = 8\pi B_2^2$ & 強 & 既に 2\%で観測一致 \\
$\Omega_{\rm dm} \sim \pi^2/36$ & 弱 & Flat 性と矛盾 \\
$\alpha^{-1}\Omega_\Lambda |B_2| = 16$ & 弱 & Post-hoc fit \\
\bottomrule
\end{tabular}
\end{center}

\subsection{Look-elsewhere effect}

多くの「偶然的一致」を報告することの危険性:

\begin{itemize}
\item $N$ 個の候補式を試し最良を選ぶと，$P(\text{偶然一致}) \sim N \cdot p_0$
\item 十分大きな $N$ では「すべての数」が何らかの arithmetic 表現と
マッチする
\item 本 note では $\sim 20$ 個の候補式を試した．偶然確率を $10\%$
とすると look-elsewhere factor $\sim 2$
\end{itemize}

これは警戒すべき．ただし:
\begin{itemize}
\item 先行研究で\textbf{観測前に}提案された式 ($\alpha^{-1}$ 公式,
$\Omega_\Lambda = 2\pi/9$) は post-hoc fitting ではない
\item 新規観察 ($9!$ identity, E$_8$, 252) は consistency checks
であり，既存仮説の強化
\item 独立な物理量 (electromagnetism, cosmology, strings) に同じ
arithmetic が現れるのは look-elsewhere で説明しきれない
\end{itemize}

\subsection{戦略}

Wright Brothers プロジェクトの健全性を保つため:
\begin{enumerate}
\item \textbf{予測の文書化}: 新しい arithmetic 表現は\textbf{観測前}に
文書化する
\item \textbf{厳密な検証}: 数値一致だけでなく理論的導出を要求する
\item \textbf{Open-source}: 全ての試行 (成功/失敗) を公開し，
selection bias を minimize
\item \textbf{Honest な reporting}: 本 note のように weak な観察も
明示する
\end{enumerate}

\section{面白い発見ランキング (主観)}

これまで見つけた observation の「面白さ」を主観的に順序付け:

\begin{enumerate}
\item \textbf{E$_8$ roots $= 240 = 8/|B_8|$} (新発見，group theoretic)
\item \textbf{SM gauge \# $= 12 = 2/|B_2|$} (深い連関の可能性)
\item \textbf{$12 \cdot 120 \cdot 252 = 9!$} (clean identity)
\item \textbf{Muon $g-2 \sim 252 = 6/|B_6|$} (先行，強い数値一致)
\item \textbf{$\Omega_\Lambda = 8\pi B_2^2$} (簡潔表現)
\item \textbf{$691$ での断絶} (suggestive だが物理対応不明)
\item \textbf{世代-$\zeta$ 対応仮説} (speculative，タウ g-2 予測)
\end{enumerate}

### E$_8$ との関係の意義

発見 1 の E$_8$ root $= 240$ は特に興味深い．E$_8$ は:
\begin{itemize}
\item 最大の exceptional Lie 群
\item Heterotic string compactification で central
\item Quasi-crystal と E$_8$ root lattice
\item M-theory と E$_{10}$, E$_{11}$ hierarchy
\end{itemize}

もし Wright Brothers framework が E$_8$ を Bernoulli arithmetic で
``predict'' できるなら，これは heterotic string theory との
unexpected な connection を示す．$240 = 8/|B_8|$ が偶然でないなら，
次は $E_8 \times E_8$ or $E_7$ などの構造も Bernoulli で記述される
はず．

### 検証可能な予測

発見群から派生する検証可能な予測:
\begin{enumerate}
\item \textbf{Tau lepton $g-2$}: Wright Brothers は $|\zeta(-9)| = 1/132$
or 関連値を予測．Belle II/LHCb で将来検証
\item \textbf{$\alpha^{-1}(m_Z)$ precision}: $127.95$ の理論値は
arithmetic 予測．FCC-ee で $0.01\%$ 精度向上
\item \textbf{Higgs self-coupling}: arithmetic 予測を探索．HL-LHC で
$50\%$, FCC で $5\%$ 精度
\end{enumerate}

\section{結論}

本 note は $B_2$ 統一仮説を深掘りし，以下の observation を報告した:

\begin{enumerate}
\item Bernoulli ladder $2n/|B_{2n}| = 12, 120, 252, 240, 132$ は
$1/|\zeta(1-2n)|$
\item 最初の 3 項の積は $9!$ に一致
\item $12$ は SM gauge boson 数と一致
\item $240$ は E$_8$ root 数と一致 (\textbf{新発見})
\item $252$ はミューオン $g-2$ 異常と一致 (先行観察の再確認)
\item $\Omega_\Lambda = 8\pi B_2^2$ の簡潔表現
\item $k = 6$ で $691$ により Bernoulli ladder の整数性が破れる
(物理的対応は不明)
\end{enumerate}

これらは\textbf{Spec$(\mathbb{Z})$ alphabet 仮説}——物理定数は
arithmetic 不変量の組合せで書かれる——を支持する．とりわけ E$_8$
との合致は新しい方向を示唆し，heterotic string との unexpected な
connection の可能性がある．

ただし，look-elsewhere effect のリスクは依然大きい．本格的な
検証には (a) loop 計算から arithmetic を rigorous に導出する，
(b) 観測前に未測定量を予測する，(c) 独立な framework から
$B_2, \zeta(1-2n)$ 等の central 性を導く，等の努力が必要．

Wright Brothers プロジェクトは引き続き「観察 → 仮説 → 検証」のサイクルを
maintain し，honest な reporting で数論的宇宙論の feasibility を
探索する．

\begin{thebibliography}{99}

\bibitem{b2-note}
Wright Brothers Research Program, ``$B_2 = 1/6$ 統一仮説: 微細構造定数と
ダークエネルギー密度の共通 arithmetic 基盤'' (2026),
\texttt{b2\_unification\_ja.tex}.

\bibitem{Kummer1850}
E.~E.~Kummer, J.\ Reine Angew.\ Math.\ \textbf{40}, 131 (1850).

\bibitem{muonNote}
Wright Brothers Research Program, early observation notes on
$252 = 6/|B_6|$ and muon $g-2$ anomaly.

\bibitem{ConnesMarcolli}
A.~Connes and M.~Marcolli, \textit{Noncommutative Geometry, Quantum
Fields and Motives} (AMS, 2008).

\bibitem{BMW2021}
S.~Borsanyi \textit{et al.} (BMW Collaboration),
Nature \textbf{593}, 51 (2021).

\bibitem{FNAL2023}
D.~P.~Aguillard \textit{et al.} (Muon $g-2$ Collaboration),
Phys.\ Rev.\ Lett.\ \textbf{131}, 161802 (2023).

\end{thebibliography}

\end{document}
